\documentclass[11pt,a4paper]{article}
\usepackage[utf8]{inputenc}
\usepackage[T1]{fontenc}
\usepackage[spanish,es-tabla]{babel}
\usepackage{lmodern,amsmath,amssymb,booktabs,array}
\usepackage[margin=2.3cm]{geometry}
\usepackage[hidelinks]{hyperref}
\setlength{\parindent}{0pt}
\setlength{\parskip}{6pt}
\newcommand{\eur}{\,\text{euros}}
\begin{document}
\begin{center}{\Large\bfseries Ejercicio 1. Problema del viajante}\par
QUIROTEC · Logística y Servicios\end{center}
\section*{Ruta óptima y distancia}
Se busca un circuito que visite exactamente una vez cada una de las veinte ciudades
y regrese a su ciudad de partida. Se utilizan exclusivamente las distancias de la matriz
del enunciado, expresadas en kilómetros. Se fija A Coruña como origen para presentar la ruta.

La solución óptima tiene una distancia total de \textbf{3.946 km}, incluido el regreso
desde León hasta A Coruña. El recorrido se detalla a continuación.

\begin{center}\small
\begin{tabular}{rllrr}\toprule
Tramo & Origen & Destino & km & Acumulado \\\midrule
1 & A Coruña & Vigo & 141 & 141 \\
2 & Vigo & Ourense & 87 & 228 \\
3 & Ourense & Salamanca & 305 & 533 \\
4 & Salamanca & Badajoz & 295 & 828 \\
5 & Badajoz & Sevilla & 207 & 1.035 \\
6 & Sevilla & Málaga & 191 & 1.226 \\
7 & Málaga & Córdoba & 155 & 1.381 \\
8 & Córdoba & Madrid & 358 & 1.739 \\
9 & Madrid & Albacete & 240 & 1.979 \\
10 & Albacete & Alicante & 163 & 2.142 \\
11 & Alicante & Valencia & 158 & 2.300 \\
12 & Valencia & Castellón & 67 & 2.367 \\
13 & Castellón & Barcelona & 269 & 2.636 \\
14 & Barcelona & Lleida & 155 & 2.791 \\
15 & Lleida & Zaragoza & 142 & 2.933 \\
16 & Zaragoza & Pamplona & 171 & 3.104 \\
17 & Pamplona & Santander & 229 & 3.333 \\
18 & Santander & Burgos & 147 & 3.480 \\
19 & Burgos & León & 180 & 3.660 \\
20 & León & A Coruña & 286 & 3.946 \\
\midrule
\multicolumn{3}{r}{\textbf{Distancia total}} & \textbf{3.946} & \\
\bottomrule\end{tabular}\end{center}

La ruta contiene las veinte ciudades sin repeticiones y tiene veinte desplazamientos.
Como la matriz es simétrica, recorrer este mismo circuito en sentido inverso da la misma
distancia. Cambiar el punto de partida del circuito tampoco altera su longitud.

El archivo \texttt{03\_TSP\_optimo.xlsx} contiene los identificadores de la ruta.
Las ciudades, los destinos, la distancia de cada tramo y la suma total se obtienen con
fórmulas, mediante búsquedas en la matriz original.

\newpage
\section*{Resolución con NEOS y verificación}
Se ha seguido el procedimiento de la guía, utilizando directamente el archivo
\texttt{TSP\_datos\_ES.mps} proporcionado para la práctica. El modelo se resolvió mediante
Gurobi en el servidor NEOS. Contiene 380 variables binarias para seleccionar los arcos,
20 variables auxiliares continuas y 401 restricciones. Exige una entrada y una salida
por ciudad e incorpora restricciones de orden para impedir circuitos parciales.

Se envió el MPS original junto con los parámetros \texttt{MIPGap=0} y
\texttt{MIPGapAbs=0}, solicitando la devolución de los archivos de solución.
El trabajo 24012137 se resolvió el 23 de octubre de 2026. Su salida registra:

\begin{center}\begin{tabular}{lr}\toprule
Dato de la ejecución & Resultado \\\midrule
Motor & Gurobi 13.0.1 \\
Estado & \texttt{Optimal solution found} \\
Mejor objetivo & 3.946 km \\
Mejor cota & 3.946 km \\
Brecha de optimalidad & 0,0000\,\% \\
Tiempo del optimizador & 0,43 segundos \\
\bottomrule\end{tabular}\end{center}

La igualdad entre el objetivo y la cota, con brecha nula, certifica la optimalidad global.
El ZIP devuelto contiene diez soluciones; \texttt{model0.sol} es la óptima, con objetivo
3.946. Las otras nueve son soluciones intermedias de mayor coste.

Para reconstruir la ruta se toman los arcos con $x_{ij}=1$. Por ejemplo,
$x_{1,19}=1$ indica A Coruña--Vigo, y $x_{19,13}=1$, Vigo--Ourense.
Siguiendo los veinte arcos se obtiene el circuito de la primera página, que se trasladó
al Excel. Se comprobó que visita las veinte ciudades, vuelve al origen y suma 3.946 km.
Las distancias de los tramos y su suma se calculan mediante fórmulas en el libro.

\section*{Comparación con las heurísticas}
El vecino más próximo comienza en A Coruña y escoge en cada paso la ciudad no visitada
más cercana. En caso de empate se toma el menor identificador de la matriz. Finalmente
se vuelve al origen. El libro \texttt{01\_TSP\_vecino.xlsx} calcula esta selección con fórmulas.

\begin{center}\begin{tabular}{lrr}\toprule
Procedimiento & Distancia (km) & Exceso sobre el óptimo \\\midrule
Orden alfabético & 12.388 & 8.442 \\
Vecino más próximo & 4.348 & 402 \\
Evolutionary: ejecución en Excel & 3.998 & 52 \\
Gurobi en NEOS & 3.946 & 0 \\\bottomrule
\end{tabular}\end{center}
El óptimo reduce el recorrido del vecino más próximo en 402 km, un 9,25\,\%.

La ejecución de Evolutionary realizada en Excel Online obtuvo 3.998 km. Solver indicó
que había convergido a la solución actual y que se satisfacían todas las restricciones.
Los controles de repeticiones y de identificadores fuera de rango dieron cero.
Esta solución mejora en 350 km al vecino más próximo y supera el óptimo en 52 km.
La convergencia del método evolutivo no constituye una prueba de optimalidad global.

\small Fuentes: matriz de distancias de la práctica, p. 1; salida de NEOS del trabajo
24012137 y archivo \texttt{model0.sol}; modelo original \texttt{TSP\_datos\_ES.mps};
ejecución de Evolutionary en Excel Online.
\end{document}
